\documentclass[12pt,a4paper]{article}
\usepackage{graphicx}
\usepackage[utf8]{inputenc}
\usepackage[T1]{fontenc}
\usepackage[polish]{babel}
\usepackage{lmodern}
\usepackage[a4paper,margin=2.5cm]{geometry}
\usepackage{amsmath,amssymb}
\usepackage[normalem]{ulem}
\usepackage{graphicx}
\usepackage{listings}
\usepackage{hyperref}

\hypersetup{
    colorlinks=true,
    urlcolor=blue
}

\lstset{
    language=Python,
    basicstyle=\ttfamily\small,
    frame=single,
    breaklines=true,
    showstringspaces=false
}

\setlength{\parskip}{0.6em}
\setlength{\parindent}{0pt}

\begin{document}

\section*{Jak dbać o zdrowy sen?}

\textbf{Zatrważająca liczba ludzi została odnotowana w badaniach jako cierpiących na bezsenność. Co możemy zrobić, aby zapewnić sobie zdrowy wypoczynek i dobre przygotowanie do nadchodzącego dnia?}

\subsection*{Jak poprawić jakość snu?}

Niektórzy uważają, że \sout{bezsenności nie da się wyleczyć ani złagodzić jej wpływu na codzienne życie}. To nieprawda. \textbf{\textit{Zdrowa wieczorna rutyna}} może poprawić jakość snu. Oto kilka prostych sposobów:

\begin{itemize}
    \item Kładź się spać o podobnej porze każdego dnia.
    \item Ogranicz korzystanie z ekranów przed snem.
    \item Zadbaj o ciszę, ciemność i odpowiednią temperaturę w sypialni.
\end{itemize}

\subsection*{Jakie są skutki braku snu?}

Do najczęstszych skutków braku snu należą:

\begin{enumerate}
    \item Zmęczenie i brak energii w ciągu dnia.
    \item Problemy z koncentracją i pamięcią.
    \item Zwiększone ryzyko chorób serca i cukrzycy.
\end{enumerate}

Jeżeli doświadczasz chronicznej bezsenności, warto skonsultować się z lekarzem lub specjalistą ds. snu, aby znaleźć odpowiednie rozwiązania i strategie poprawy jakości snu. Zwróć uwagę na:

\begin{itemize}
    \item[$\square$] Problem z zasypianiem lub utrzymaniem snu.
    \item[$\square$] Częste budzenie się w nocy.
\end{itemize}

Warto skonsultować się z lekarzem, jeśli problemy te utrzymują się przez dłuższy czas i wpływają na codzienne funkcjonowanie. Pamiętaj, że zdrowy sen jest kluczowy dla ogólnego zdrowia i dobrego samopoczucia.

\subsection*{Bezsenność w Kanadzie w latach 2018 i 2020}

W badaniu kohortowym porównano objawy bezsenności przed pandemią i podczas jej pierwszej fali.

\begin{center}
\begin{tabular}{|c|r|r|}
    \hline
    Rok & Objawy bezsenności & Zespół bezsenności \\
    \hline
    2018 & 25,4\% & 16,8\% \\
    2020 & 32,2\% & 19,0\% \\
    \hline
\end{tabular}
\end{center}

\begin{figure}[h]
    \centering
    \includegraphics[width=0.8\textwidth]{wykres.png}
    \caption{Wykres przedstawiający wpływ snu na zdrowie}
\end{figure}

\textit{Czasami ludzie z bezsennością zgłaszają sny związane z rzutem ukośnym (diagram wyżej).}

Jeżeli odczuwasz problemy ze snem, możesz podzielić się swoimi doświadczeniami i strategiami radzenia sobie z bezsennością w \href{https://colab.research.google.com/}{notebooku Google Colab}. Wspólne dzielenie się wiedzą może pomóc w znalezieniu skutecznych metod poprawy jakości snu.

Z naszych doświadczeń wynika, że osoby cierpiące na bezsenność to w części matematycy specjalizujący się w kilku specyficznych wzorach, takich jak
\(\sum_{n=1}^{\infty} \frac{1}{n^2}\),
\(\int_0^\infty e^{-x^2}\,dx\)
czy
\(\lim_{x \to 0} \frac{\sin(x)}{x}\).
Zgłaszano również przypadki pracy nad takimi wzorami:

\[
\sum_{n=31}^{\infty} \frac{3}{n^2}
\]

\[
\int_0^\infty e^{-x^3}\,dx
\]

\[
\int_0^\infty \frac{\sin(x)}{x}\,dx
\]

Warto zauważyć, że te wzory mogą prowadzić do stresu i trudności w zasypianiu.

Zupełnie poważnie mówiąc, ten artykuł/zadanie, mimo że nosi znamiona opowiadania lub pracy absurdalnej, w dużej mierze zostało bardzo dobrze i rzetelnie uzupełnione dzięki podpowiedziom dostępnego w Visual Studio Code chata, co wywołuje u mnie niemałe przerażenie. Uważam, że zabrnąłem już w absurd za daleko, dlatego na sam koniec programuję ostatnią część tego zadania.

\begin{lstlisting}
x = int(input("Podaj liczbę: "))
y = x**3
print(f"Wynik potęgowania: {y}")
print("Miłego dnia i zdrowego snu!")
\end{lstlisting}

\end{document}